\section{Introduction}
\label{sec:intro}

A vision--language model that answers ``is there a mouse in the image?'' wrongly is more
useful if it says so. Selective prediction does this: a score ranks the model's answers, a
threshold is calibrated on held-out data, and the model answers only above it. Conformal and
risk-controlling variants \citep{bates2021rcps,angelopoulos2021ltt,angelopoulos2024crc} certify,
with a finite-sample guarantee that makes no assumption on the score or the model, that the
error rate among the answers that are emitted stays below a target $\alpha$.

The guarantee is paid for in coverage, and part of the bill cannot be reduced. Kotte
\citep{crcimposs} shows that when the base error rate $\mu$ exceeds $\alpha$, any predictor that
may only emit or abstain the model's own answer must abstain on at least
$(\mu-\alpha)/(1-\alpha)$ of inputs, whatever the score, the concentration inequality or the
calibration budget (Proposition~\ref{prop:floor} below). The proof uses one premise: an
abstained item is an item whose model answer stays wrong. A predictor that may emit a
\emph{different} answer is outside the premise.

This paper studies that possibility in the simplest setting where it can be made exact: yes/no
object-existence questions, with an open-vocabulary detector as a second, distinct source of
evidence. \emph{Certified-Correction Risk Control} (CCRC) accepts the model's answer where a
correctness score is high, replaces it with the detector's answer where the detector's evidence
is decisive, abstains otherwise, and certifies the error rate of the whole emitted set, accepted
and substituted answers together.

\paragraph{What the paper contributes}
Our aim is a result that survives a strict reading of the protocol. Four design decisions follow
from that aim, and each is a contribution in its own right.
\begin{enumerate}
\item \emph{The proved procedure is the evaluated procedure.} Theorem~\ref{thm:validity} covers
  policy families indexed by threshold \emph{values} fixed before calibration under i.i.d.
  sampling of calibration units. Every reported number uses such a family: thresholds and repair
  gate are read off the fit fold, never off the calibration sample whose errors are counted.
\item \emph{The sampling unit is stated and respected.} POPE contributes six questions per
  image; folds here are unions of whole images, we measure the within-image dependence of
  errors, and a known-population simulation shows how clustering would degrade the guarantee and
  how a one-item-per-image calibration restores it.
\item \emph{The repair gate is not tuned on the evaluation data.} We report three ways of
  fixing it that keep the guarantee: a constant chosen on a different benchmark, a choice made
  on the fit fold, and a union bound over a pre-specified set of gates.
\item \emph{Uncertainty is stated at the level of the data.} Gains are accompanied by an outer
  image-level bootstrap, because hundreds of random splits of one benchmark measure the
  procedure's sensitivity to the split, not the benchmark's sampling error.
\end{enumerate}

\paragraph{What we find}
\begin{itemize}
\item With the gate fixed on an external benchmark at the target's calibration size, CCRC certifies $+0.8$ to $+7.3$ points
  more coverage than filtering in all twelve (cell, $\alpha$) combinations, and the image-bootstrap interval excludes zero in
  seven. The gain is about one point on cells with $150$--$200$ calibration items and up to $+7.3$ points with $500$ or more
  (\S\ref{sec:main}).
\item Choosing the gate on the fit fold, or certifying three gates by a union bound, is valid and helps with $500$ or more calibration
  items ($+2.4$ to $+8.8$ and $+2.0$ to $+4.4$ points) but costs more than it returns with $150$--$200$
  ($-6.2$ to $+9.4$ and $-16.5$ to $+4.4$ points).
\item Much of the repair is confirmation: $22$--$100\%$ of repaired items are items on which the detector agrees with the model.
  With a small gate the answer is changed for under $0.5\%$ of test items; with a large gate for $2.4$--$5.7\%$, of which $81$--$93\%$
  are right, and $4$--$5\%$ of the model's false-positive errors are corrected (\S\ref{sec:accounting}).
\item The start of the fixed sequence matters more than the repair on small folds. Starting where a single error can be
  tolerated ($k\ge38$ at $\alpha=\delta=0.10$) rather than where none can ($k\ge22$) raises filtering coverage on the smallest cell
  from $22\%$ to $41\%$ and removes an apparent loss of repair on AMBER, where the conventional start makes CCRC abort on $72\%$
  of splits (\S\ref{sec:start}).
\item Certified coverage comes mostly from the score: a detection-grounded score certifies $68\%$ on POPE-1500 against $41\%$ for
  model confidence and $0.5\%$ for the image--text similarity score of ConfLVLM, while certified repair adds $1.4$ to $8.8$ points
  (\S\ref{sec:score}).
\item In free-form generation, repairing a claim mid-caption increases downstream hallucination in a matched-prefix experiment
  (\S\ref{sec:sequential}), which blocks a fixed-point argument for a sequential guarantee.
\end{itemize}

\paragraph{Scope}
Everything here concerns binary existence questions, one detector, one family of backbones and
benchmarks whose images overlap ($1{,}457$ distinct images across all cells: the $500$ POPE images, nested across the POPE cells, and $957$ AMBER images). We do not claim that correction is
generally useful, only that under a protocol whose guarantee we can state it is valid and its
gain is measured.
